\chapter{Kilometre Line and Stationing}

\section{Motivation}

\begin{remark}
Railway infrastructure is not only described geometrically, but also by
stationing.
\end{remark}

\begin{remark}
The kilometre line defines the longitudinal reference used for
addressing positions along a railway route.
\end{remark}

\begin{remark}
It must not be confused with the physical track centerline.
\end{remark}

\section{Definition}

\begin{definition}[Kilometre line]
A kilometre line is a longitudinal reference alignment that assigns a
station value to positions along a railway route.
\end{definition}

\begin{remark}
The kilometre line may coincide with a track centerline, but this is not
required.
\end{remark}

\section{Relation to Track Alignments}

\begin{remark}
Several track alignments may refer to the same kilometre line.
\end{remark}

\begin{remark}
This is common in double-track lines, junctions, station areas, and
switching zones.
\end{remark}

\begin{remark}
In single-track cases, the kilometre line and the track centerline may
be identical.
\end{remark}

\section{Outer Stationing}

\begin{remark}
When the physical track alignment differs from the kilometre line,
positions along the track must be related back to the kilometre
reference.
\end{remark}

\begin{definition}[Outer stationing]
Outer stationing describes the mapping between a track alignment and an
external kilometre reference.
\end{definition}

\section{Topological Role}

\begin{remark}
The kilometre line acts as a topological and administrative reference
rather than merely a geometric object.
\end{remark}

\begin{remark}
It provides continuity across track changes, parallel tracks, and
operationally relevant segments.
\end{remark}

\section{Implication for AIM}

\begin{summary}
Within AIM, the kilometre line is treated as a reference structure
separate from physical track alignments.

This distinction is essential for modelling multi-track situations,
switches, stationing jumps, and real-world railway project data.
\end{summary}

\chapter{Kilometre Line, Stationing, and Spatial Reference}

\section{Motivation}

\begin{remark}
Railway infrastructure is not only described by geometry, but also by
stationing.
\end{remark}

\begin{remark}
A point on a railway route is usually not addressed by global
coordinates, but by a kilometre value along a reference line.
\end{remark}

\begin{remark}
Thus, railway modelling must distinguish between:
\begin{itemize}
\item geometric location,
\item spatial reference system,
\item and operational addressing.
\end{itemize}
\end{remark}

\section{Three Reference Layers}

\begin{definition}[Spatial reference layer]
The spatial reference layer defines where geometry lives.
It may be based on:
\begin{itemize}
\item a planar projection,
\item a local engineering coordinate system,
\item or, in a long-term view, a reference ellipsoid.
\end{itemize}
\end{definition}

\begin{definition}[Geometric alignment layer]
The geometric alignment layer describes the physical or designed track
geometry, typically through curvature, pose, profile, and cant.
\end{definition}

\begin{definition}[Stationing layer]
The stationing layer assigns address values to positions along a railway
route.
\end{definition}

\begin{summary}
Spatial reference answers: where is it?  
Geometry answers: what shape is it?  
Stationing answers: how is it addressed?
\end{summary}

\section{Kilometre Line}

\begin{definition}[Kilometre line]
A kilometre line is a longitudinal reference structure used to assign
station values to positions within a railway route or project.
\end{definition}

\begin{remark}
The kilometre line may coincide with a physical track centerline, but
this is not required.
\end{remark}

\begin{remark}
It is therefore incorrect to treat the kilometre line simply as another
track alignment.
\end{remark}

\section{Relation to Track Geometry}

\begin{remark}
Several physical track alignments may refer to the same kilometre line.
\end{remark}

\begin{remark}
This occurs in:
\begin{itemize}
\item double-track lines,
\item station areas,
\item junctions,
\item switches and crossings,
\item temporary or construction-stage alignments.
\end{itemize}
\end{remark}

\begin{remark}
Conversely, in simple single-track situations, the kilometre line and
track centerline may be identical.
\end{remark}

\section{Outer Stationing}

\begin{definition}[Outer stationing]
Outer stationing describes the mapping between a physical track
alignment and an external kilometre reference.
\end{definition}

\begin{remark}
Outer stationing is required whenever the track geometry and the
kilometre reference do not coincide.
\end{remark}

\begin{remark}
In such cases, a point on the track has at least two relevant
longitudinal coordinates:
\begin{itemize}
\item its intrinsic track coordinate \(s_{\mathrm{track}}\),
\item its assigned kilometre value \(s_{\mathrm{km}}\).
\end{itemize}
\end{remark}

\section{Relation to Ellipsoid and CRS}

\begin{remark}
The kilometre line is independent of the coordinate reference system in
which the geometry is represented.
\end{remark}

\begin{remark}
A kilometre value is not a projected coordinate.
It is an addressing value along a railway-specific reference structure.
\end{remark}

\begin{remark}
Thus, even if geometry is formulated on a reference ellipsoid, in a
Gauss--Krüger projection, or in a local engineering CRS, the kilometre
line remains a separate semantic layer.
\end{remark}

\section{Reference Separation Principle}

\begin{proposition}
A railway project requires separation of at least three reference
concepts:
\[
\text{spatial reference}
\quad\neq\quad
\text{geometric alignment}
\quad\neq\quad
\text{stationing reference}.
\]
\end{proposition}

\begin{remark}
Confusing these layers leads to systematic modelling errors, especially
in multi-track and junction situations.
\end{remark}

\section{Implication for AIM}

\begin{remark}
Within AIM, the kilometre line is treated as a reference object, not as
the primary physical track geometry.
\end{remark}

\begin{remark}
Track alignments may be assigned to kilometre lines by explicit
stationing relations.
\end{remark}

\begin{remark}
This allows AIM to model cases where:
\begin{itemize}
\item one kilometre line serves multiple tracks,
\item a track has no independent kilometre line,
\item stationing jumps occur,
\item geometric and operational continuity differ.
\end{itemize}
\end{remark}

\section{Relation to the Seven-Line Model}

\begin{remark}
The kilometre line plays a central role in the seven-line import and
sorting model.
\end{remark}

\begin{remark}
A simplified structure is:
\[
\text{km-line}
\rightarrow
\begin{cases}
\text{right track} \rightarrow (\text{profile}, \text{cant}),\\
\text{left track}  \rightarrow (\text{profile}, \text{cant}),\\
\text{outer stationing}.
\end{cases}
\]
\end{remark}

\begin{remark}
This expresses that the kilometre line provides the addressing context,
while the tracks provide the physical geometry and engineering state.
\end{remark}

\section{Key Insight}

\begin{proposition}
The kilometre line is not primarily a geometric curve, but an addressing
structure for railway infrastructure.
\end{proposition}

\begin{summary}
The km-line, the geometric alignment, and the spatial reference system
serve different roles.

The km-line addresses railway infrastructure.

The alignment describes physical geometry.

The CRS or ellipsoid defines the spatial embedding.

Keeping these layers separate is essential for robust railway modelling,
especially in multi-track, station, and junction contexts.
\end{summary}
